\chapter{Experimental Setup, Data, and Evaluation Protocol}
\label{ch:experimental-setup}

What did the FLFM cameras record, and which of those measurements can be used
to assess the reconstruction in Chapter~5? Answering that requires more than
the reconstructed image. We need the optical layout, the acquisition record,
and the processing run that produced each result. This chapter describes those
records and the comparisons they allow.

\section{Optical configuration}
\label{sec:optical-configuration}

The instrument has one illumination path and three camera paths. A Coherent
SAPPHIRE 488-200 CDRH laser supplies $488\,\mathrm{nm}$ light. An expanding
lens and a $40\times$ objective focus it onto a thin sample held between two
coverslips. That objective also collects light for the main Fourier
light-field microscopy (FLFM) camera or, when the flip mirror is engaged, the
ThorCam. A second $40\times$ objective faces the sample from the other side
and leads to the Validation Camera. The component inventory gives
$20\,\mathrm{mW}$ as a commonly used laser setting. It does not give the
power at the sample or the setting for individual acquisitions.

In the main collection path, a two-lens relay images the objective's Fourier
(pupil) plane onto the microlens array. The installed Advanced Microoptics
Systems GmbH APO-Q-P3000-R23.5 has $10\times10$ fused-silica microlenses on
a square grid with $3\,\mathrm{mm}$ pitch. Its working focal length is
$55\,\mathrm{mm}$; the flat side faces away from the incoming collimated
light. Another two-lens relay takes the array output to the Teledyne-FLIR
Blackfly S BFS-PGE-244S8M main camera. In the ray-optics design, the lenses
before the array have focal lengths of $250\,\mathrm{mm}$ and
$750\,\mathrm{mm}$, and those after it have focal lengths of
$400\,\mathrm{mm}$ and $150\,\mathrm{mm}$.
Individual lens part numbers and measured separations are still missing.

Before the microlens array, a flip-mounted mirror can send the collected light
to the ThorCam branch. This branch avoids the array and focuses through a
regular lens onto a Thorlabs CS2100M-USB camera. The ray-optics script models
that lens as $300\,\mathrm{mm}$. The objective and first relay are shared
with the FLFM path; everything after the mirror belongs to one branch or the
other. Switching the mirror means that the two cameras capture at different
times. Across the sample, the validation path passes through its own
$40\times$ objective and a lens modelled as $100\,\mathrm{mm}$ before
reaching a Teledyne-FLIR Blackfly S BFS-U3-122S6M camera.

The objective used for both the illumination/collection path and the validation camera path is ThorLabs N40X-PF $40\times$ objective: focal length
$5\,\mathrm{mm}$, numerical aperture 0.75, working distance
$0.66\,\mathrm{mm}$, and entry pupil $7.5\,\mathrm{mm}$. The
available records also leave open how excitation enters the shared objective,
which filters are installed, and what contrast the Validation Camera records.
The ray-optics scripts establish the path order, but they cannot settle these
hardware details or the spacing of the assembled lenses.

\section{Samples, acquisitions, and view extraction}
\label{sec:sample-acquisition}

The acquisition archive contains fluorescent-bead scenes. Filenames often
identify bead size, density, medium, and exposure; many specify a nominal
$2\,\mathrm{\mu m}$ bead size. Among the candidate families are scenes
labelled \emph{eyedrop} and scenes acquired with chicken-breast material.
These labels help locate an image, but they are not a sample-preparation
record: bead concentration, medium thickness, coverslip thickness, and
delivered illumination remain unknown for a selected capture. The physical
sample sits between two coverslips. A simple simulation includes glass on
both sides, with assumed thickness and refractive index; those inputs are not
measurements of the coverslips used for acquisition.

Acquisitions are grouped by date. A date folder can contain FLFM images in
\texttt{main}, conventional captures in \texttt{thorcam}, and Validation Camera
images in \texttt{validation} or \texttt{val}. Extracted views, rectification
runs, and analysis products appear under \texttt{algo\_run}. Folder proximity
alone does not establish that two files depict the same capture. A comparison
needs the scene identity and acquisition date, especially for ThorCam images,
which are taken after the flip mirror is moved.

Before reconstruction, the FLFM sensor image is cut into a square grid of
elemental views. Write a view as $V_{ij}(y,x)$, where
$i,j\in\{0,\ldots,N-1\}$ index the grid and $(y,x)$ are pixel coordinates
within a view. The resulting array has shape $(N,N,H,W)$; a selected exposure
stack adds an exposure axis. The physical $10\times10$ array alone does not
fix the extracted grid for any particular run. To reproduce that step, we
need its center, pitch or patch geometry, orientation, and excluded views.
We also need to know which intensity image was supplied. An unmodified sensor
image, a dark-subtracted image, and a display-normalized image carry different
values. The archive contains dark-frame and gain-calibration files, but their
use has to be checked for each dataset. The same applies to exposure fusion:
it belongs to a run only when its saved configuration enables it.

\section{Reconstruction and selection records}
\label{sec:reconstruction-records}

Chapter~5 describes one reconstruction sequence: global translation
alignment, then local pixelwise refinement and averaging. A run may save the
average after global alignment as well as the final average after local
refinement. Comparing them shows what the local step changed, provided both
come from the same run and input views. A reported image therefore needs a
source capture, an extracted view grid, a processing configuration, and the
saved numerical output.

The run archive may contain a sweep specification, an experiment manifest,
and generation and candidate identifiers. Individual runs may also have
\texttt{metric\_record.json}, \texttt{config.json}, and \texttt{workspace.npz}.
The configuration gives requested settings, while effective values also depend
on the implementation and input geometry. The workspace and numerical view
stack show which data were averaged. Preview images help with visual
inspection, but may have a different intensity scale. To use a run as thesis
evidence, its record needs the code revision, effective parameters,
preprocessing, and both intermediate and final outputs. Some older products
are missing a source revision or a complete history of figure preparation.
A place on a sweep leaderboard cannot fill those gaps.

Parameter sweeps and image-quality scores narrow the set of images worth
inspecting. A ranking applies only to the candidates that completed for that
dataset. Sharpness or contrast alone cannot tell us whether the image
reproduces the bead scene faithfully. For each example in the results chapter,
the candidate set, ranking rule, failed or missing runs, and any manual choice
need to be reported. If an external-camera image or bead annotation helped
select the candidate, comparing it with that same reference is part of the
selection process, not an independent test.

\section{Evaluation protocol}
\label{sec:evaluation-protocol}

The most direct comparison stays within one FLFM run: the average after global
alignment against the final rectified image. They share the acquisition and
parameter record, so the comparison shows the contribution of local
refinement. A visual panel needs a stated crop and display scale. A numerical
comparison needs the saved images, a defined metric and direction, and the
same stated region in both images. Local remapping can leave zero-filled
boundaries; including different amounts of those boundaries would change the
score for reasons unrelated to the refinement itself.

The ThorCam and Validation Camera offer two other views of the sample. ThorCam
uses a separate branch of the shared collection path and is captured at a
different time. The Validation Camera looks through the opposite objective.
Either image may differ from the FLFM reconstruction in pixel sampling,
orientation, field of view, and intensity scale. A comparison therefore needs
documented resampling, geometric registration, intensity handling, and a
common crop or validity mask. Evaluation registration may use scale and
rotation; the global reconstruction stage itself estimates translations only.
Saved validation workspaces may contain transformations and overlays, but the
registration and metric must be read from the selected comparison rather than
inferred from current software defaults. Sample change between the sequential
FLFM and ThorCam captures is another limit on interpretation.

The optical paths and archive structure are known well enough to locate
candidate comparisons. The final acquisition table still needs sample
preparations, exposure and power settings, extraction geometry, camera
registration, and evaluation metrics for the selected examples. Those records
will determine which comparisons in the next chapter can support a result.
